\documentclass[12pt]{article}
\usepackage[utf8]{inputenc}
\usepackage[T1]{fontenc}
\usepackage[margin=1in]{geometry}
\usepackage{setspace}
\usepackage{xcolor}
\usepackage[colorlinks=true, urlcolor=blue!60!black, linkcolor=blue!60!black]{hyperref}
\setlength{\parindent}{2em}
\setlength{\parskip}{1.25em}
\renewcommand{\baselinestretch}{1.0}
\begin{document}
\begin{center}
{\LARGE \textit{United States v. Comstock}, 560 U.S. 126 (2010)\par}
\vspace{0.8em}
{\large Case Brief\par}
\end{center}
\vspace{1em}

\section{Memory Jogger}
Federal prisoners nearing release challenged a statute allowing their indefinite civil commitment as sexually dangerous; the Court held 7--2 that the Necessary and Proper Clause authorized it as a rational, narrow incident of Congress's power to run federal prisons and act as a responsible custodian (\href{https://supreme.justia.com/cases/federal/us/560/126/}{Justia}).

\section{Detailed Case Facts}
In 2006, Congress enacted 18 U.S.C. \S~4248 as part of the Adam Walsh Child Protection and Safety Act, authorizing a federal district court to order the indefinite civil commitment of a ``sexually dangerous'' person already in the custody of the Federal Bureau of Prisons (BOP)---including federal prisoners nearing the expiration of their criminal sentences, defendants found mentally incompetent to stand trial, and persons whose federal criminal charges were dismissed solely due to mental condition (\href{https://www.law.cornell.edu/supct/html/08-1224.ZO.html}{560 U.S. at 129--31}). To obtain commitment, the Department of Justice must certify that the individual has engaged or attempted to engage in sexually violent conduct or child molestation, suffers from a serious mental illness, abnormality, or disorder, and as a result would have serious difficulty refraining from such conduct if released. The government must prove those elements by clear and convincing evidence at a federal hearing. Under \S~4248(d), the Attorney General must make all reasonable efforts to transfer custody to the prisoner's state of domicile or trial; federal confinement continues only if no state assumes responsibility and until the detainee is no longer sexually dangerous.

\section{Procedural History}
Acting under \S~4248, the federal government instituted civil-commitment proceedings in the United States District Court for the Eastern District of North Carolina against five federal detainees---Graydon Earl Comstock, Jr., and four others---whose federal prison terms were about to expire or whose charges had been dismissed due to incompetency (\href{https://www.law.cornell.edu/supct/html/08-1224.ZO.html}{560 U.S. at 131--32}). Each respondent moved to dismiss, arguing that \S~4248 exceeded Congress's enumerated powers under Article I and violated due process and equal protection. The District Court agreed that Congress lacked Article I authority to enact \S~4248 and dismissed the proceedings. The Fourth Circuit affirmed solely on the Article I ground (551 F.3d 274 (2009)), without reaching the due process or equal protection claims. The Supreme Court granted certiorari limited to Congress's power under Article I.

\section{Judicial Votes}
The Supreme Court \textbf{reversed and remanded} by a \textbf{7--2} vote. The case was argued on January 12, 2010, and decided on May 17, 2010 (\href{https://supreme.justia.com/cases/federal/us/560/126/}{Justia}; \href{https://www.law.cornell.edu/supremecourt/text/08-1224}{Cornell LII}):

\begin{itemize}
    \item \textbf{Majority Opinion:} Justice Breyer delivered the opinion of the Court, joined by Chief Justice Roberts and Justices Stevens, Ginsburg, and Sotomayor.
    \item \textbf{Concurring Opinions:} Justice Kennedy filed an opinion concurring in the judgment. Justice Alito filed an opinion concurring in the judgment.
    \item \textbf{Dissenting Opinion:} Justice Thomas filed a dissenting opinion, joined by Justice Scalia in all but Part III--A--1--b.
\end{itemize}

\section{Holding}
\textbf{Questions Presented:}

Does the Necessary and Proper Clause (U.S. Const. art. I, \S~8, cl. 18) grant Congress authority to enact 18 U.S.C. \S~4248, authorizing the federal civil commitment of a mentally ill, sexually dangerous federal prisoner beyond the date the prisoner's criminal sentence would otherwise expire?

\textbf{Holding:} \textbf{Yes.} Relying on five cumulative considerations, the Court held that the Necessary and Proper Clause empowers Congress to enact \S~4248 as ``a means that is rationally related to the implementation of a constitutionally enumerated power'' (\href{https://supreme.justia.com/cases/federal/us/560/126/}{560 U.S. at 134, 149--50}). The Court reversed the Fourth Circuit and remanded for consideration of respondents' preserved due process and equal protection claims.

\section{Analysis of Opinions}
\subsection{Majority Opinion (Justice Breyer)}

Writing for a five-Justice majority, Justice Breyer grounded \S~4248 in a five-factor structural inquiry rooted in \textit{McCulloch v. Maryland}, 17 U.S. (4 Wheat.) 316 (1819), and \textit{Sabri v. United States}, 541 U.S. 600 (2004):

\begin{enumerate}
  \item \textbf{Breadth of the Necessary and Proper Clause:} Under \textit{McCulloch}, the Clause grants Congress broad authority to enact laws ``convenient, or useful'' or ``conducive'' to the beneficial exercise of enumerated powers---tested by whether the statute is ``rationally related to the implementation of a constitutionally enumerated power'' (\href{https://supreme.justia.com/cases/federal/us/560/126/}{560 U.S. at 133--35}). Although the Constitution does not expressly empower Congress to create most federal crimes or build prisons, Congress indisputably possesses implied authority to criminalize violations of enumerated powers, erect federal prisons, and regulate prison administration and safety.
  \item \textbf{Longstanding History of Federal Involvement:} Since 1855---and continuously since the 1949 statute upheld in \textit{Greenwood v. United States}, 350 U.S. 366 (1956), and its 1984 revision---Congress has provided mental-health care and civil commitment for mentally ill federal prisoners whose release would endanger others (\href{https://www.law.cornell.edu/supct/html/08-1224.ZO.html}{560 U.S. at 137--42}). Section~4248 was a modest addition to this settled statutory framework.
  \item \textbf{Sound Reasons as a Federal Custodian:} Having lawfully taken custody of a prisoner, the Federal Government acts as a responsible custodian when it prevents the release of a mentally ill, sexually dangerous inmate who threatens nearby and wider communities (\href{https://supreme.justia.com/cases/federal/us/560/126/}{560 U.S. at 142--43}).
  \item \textbf{Accommodation of State Interests:} Rejecting respondents' Tenth Amendment objection, Breyer noted that powers delegated by the Necessary and Proper Clause are not ``reserved to the States,'' and \S~4248(d) explicitly defers to state sovereignty by requiring the Attorney General to inform relevant states and transfer custody whenever a state agrees to assume responsibility (\href{https://supreme.justia.com/cases/federal/us/560/126/}{560 U.S. at 143--45}).
  \item \textbf{Narrow Scope and Attenuation Limit:} The Court rejected a rigid rule forbidding Congress from acting more than one implied step away from an enumerated power. Because \S~4248 applies only to the small class of individuals already in federal custody, the links to enumerated powers are not too attenuated and do not confer a general federal police power (\href{https://www.law.cornell.edu/supremecourt/text/08-1224}{560 U.S. at 146--49}).
\end{enumerate}

\subsection{Opinion Concurring in the Judgment (Justice Kennedy)}

Justice Kennedy agreed that \S~4248 is constitutional as a narrow, discrete exercise of federal custodial responsibility that respects state primacy through \S~4248(d) (\href{https://www.law.cornell.edu/supct/html/08-1224.ZC.html}{560 U.S. at 150--55}). He declined to join Breyer's opinion, however, warning that importing ``rationally related'' or ``rational basis'' rhetoric from Due Process or Equal Protection review risks obscuring the distinct structural limits of federalism. Under the Necessary and Proper Clause, the strength of the link to an enumerated power---not mere logical possibility---is essential.

\subsection{Opinion Concurring in the Judgment (Justice Alito)}

Echoing Kennedy's unease with Breyer's broad phrasing, Justice Alito emphasized that ``necessary'' requires a ``substantial link'' rather than a merely conceivable rational basis (\href{https://www.law.cornell.edu/supct/html/08-1224.ZC1.html}{560 U.S. at 155--58}). Alito found that link satisfied because federal imprisonment itself frequently severs an inmate's ties to any single state, creating a practical gap in state civil-commitment jurisdiction that the Federal Government itself helped produce.

\subsection{Dissenting Opinion (Justice Thomas)}

Justice Thomas argued that \S~4248 fails \textit{McCulloch}'s two-part test because it is neither ``plainly adapted'' to executing an identifiable enumerated power nor ``consistent with the letter and spirit of the Constitution'' (\href{https://www.law.cornell.edu/supct/html/08-1224.ZD.html}{560 U.S. at 158--80}). Criticizing Breyer's multi-factor balancing test for never identifying a single enumerated power that \S~4248 executes, Thomas contended that once a federal prisoner completes his sentence, the enumerated power justifying his conviction is exhausted; civil commitment of sexually dangerous citizens is a quintessentially local police power reserved to the States. (Justice Scalia declined to join Part III--A--1--b, which expressed Thomas's stricter originalist view of the Commerce Clause.)

\subsection{Subsequent Developments}

\begin{itemize}

\item \textbf{\textit{Comstock} Distinguished (\textit{NFIB v.\ Sebelius}):} Two years later, in \href{https://supreme.justia.com/cases/federal/us/567/519/}{\textit{National Federation of Independent Business v. Sebelius}}, 567 U.S. 519 (2012), Chief Justice Roberts---who had joined Breyer's majority in \textit{Comstock}---joined by the four joint dissenters on this point, rejected the federal government's attempt to sustain the Affordable Care Act's individual health-insurance mandate under the Necessary and Proper Clause. Roberts sharply distinguished \textit{Comstock} as an exercise of authority ``derivative of, and in service to, a granted power,'' holding that even if the mandate was ``necessary,'' compelling inactive individuals to enter commerce to create the predicate for regulation was not ``proper.''

\item \textbf{\textit{Comstock} Applied (\textit{United States v.\ Kebodeaux}):} One year later, however, the Court reaffirmed \textit{Comstock}'s core framework in \href{https://supreme.justia.com/cases/federal/us/570/387/}{\textit{United States v. Kebodeaux}}, 570 U.S. 387 (2013), holding 7--2 (again via Justice Breyer) that the Military Regulation Clause combined with the Necessary and Proper Clause empowered Congress to apply federal sex-offender registration requirements (SORNA) to a former Air Force serviceman who had already completed his court-martial sentence, because he remained subject to a preexisting federal registration obligation. Chief Justice Roberts and Justice Alito concurred in the judgment to stress that Congress lacks a general federal police power.

\end{itemize}

\section{Examples: Future Applications}
\begin{enumerate}

\item \textbf{Same-Side (Within Congress's Power): Commitment of Violently Dangerous Federal Prisoners.} Congress amends federal law to allow civil commitment, at the end of a federal sentence, of any prisoner in Bureau of Prisons custody who suffers from a serious mental illness and, as a result, would pose a serious danger of violent conduct if released. As in \S~4248, the government must prove dangerousness by clear and convincing evidence, and the Attorney General must transfer custody to any willing State. Every \textit{Comstock} consideration points toward validity: the statute rationally serves the federal power to imprison and care for federal prisoners, extends a long tradition of federal commitment of dangerous mentally ill prisoners, reflects the government's role as a responsible custodian, accommodates state interests, and is confined to people already in federal custody.

\item \textbf{Same-Side (Within Congress's Power): Holding a Contagious Prisoner Past the Release Date.} A federal inmate due for release next week develops a highly contagious, dangerous infectious disease. A federal statute allows the Bureau of Prisons, after a hearing, to keep such an inmate in medical isolation beyond the end of his sentence until he is no longer contagious, unless the health authorities of his home State agree to take custody. Under \textit{Comstock}, Congress may act as a responsible custodian of those it lawfully holds, protecting the communities into which a prisoner would be released. The measure is narrow, limited to people already in federal custody, temporary, and deferential to state health authorities, so it is a necessary and proper incident of running federal prisons rather than a general police power.

\item \textbf{Opposite-Side (Beyond Congress's Power): Commitment of State Felons via a Firearm Nexus.} Congress enacts a statute authorizing federal courts to impose lifetime federal civil commitment on any person who has ever been convicted of a \emph{state} violent felony if the person once purchased a firearm that traveled in interstate commerce, even if the person was never in federal custody and the state declines to commit him. A challenger attacks the statute as exceeding Article I. Under \textit{Comstock}, \textit{Kebodeaux}, and \textit{NFIB}, the statute is unconstitutional. Unlike \S~4248 in \textit{Comstock}, the law is not tied to existing federal custody or responsible management of federal prisons, lacks any historical pedigree of federal supervision over state-only offenders, intrudes directly into the states' traditional police power, and uses an attenuated past commercial nexus to manufacture a sweeping federal civil-commitment power that is neither ``substantially linked'' (Alito/Kennedy) nor ``proper'' (\textit{NFIB}).

\item \textbf{Opposite-Side (Beyond Congress's Power): Commitment of People Never Charged.} Congress authorizes the Attorney General to seek federal civil commitment of any person whom the FBI investigated for a federal sex offense but decided not to charge, if a court finds the person sexually dangerous. The people covered were never convicted, never held by the Bureau of Prisons, and are not subject to any continuing federal obligation. \textit{Comstock} stressed that \S~4248 reached only a small class already in federal custody, which kept its link to enumerated powers from becoming too attenuated. Without that custodial link, the statute does not carry any enumerated power into execution; civil commitment of dangerous persons at large is the kind of general police power the Court said remains with the States.

\item \textbf{Fence-Sitter (Close Call): New Registration Duties for Long-Released Federal Offenders.} Congress requires every person ever convicted of a federal sex offense to register with federal authorities, including people who completed their sentences decades ago and were released without conditions, before any federal registration law existed. In \textit{Kebodeaux}, the Court upheld applying registration rules to a former serviceman because he remained subject to a preexisting federal registration obligation when released. Here, no such continuing obligation exists, so the link to the power that supported the original conviction is weaker. On the other hand, the statute addresses the consequences of federal convictions, is less intrusive than commitment, and builds on a federal registration tradition. Whether the link is close enough is genuinely uncertain.

\end{enumerate}

\section{Critique}
\textit{Comstock} exposes a deep methodological fault line in structural constitutional law. Breyer's five-factor ``taken together'' test is pragmatic and flexible, but critics echo Justice Thomas's objection: if Congress can stack implied power upon implied power (enumerated power $\rightarrow$ federal crime $\rightarrow$ federal prison $\rightarrow$ prison medical care $\rightarrow$ post-sentence civil commitment), what principled stopping point remains? By contrast, Kennedy and Alito's concurrences---and Roberts's later opinion in \textit{NFIB}---recast ``proper'' as an independent federalism check preventing Congress from converting incidental custodial powers into a general police power.

\section{Key Quotations}
\begin{quote}
``Taken together, these considerations lead us to conclude that the statute is a `necessary and proper' means of exercising the federal authority that permits Congress to create federal criminal laws, to punish their violation, to imprison violators, to provide appropriately for those imprisoned, and to maintain the security of those who are not imprisoned but who may be affected by the federal imprisonment of others.''
\end{quote}
\noindent --- Justice Breyer, for the Court (\href{https://www.law.cornell.edu/supct/html/08-1224.ZO.html}{560 U.S. at 149}).

\begin{quote}
``Concluding that a relation can be put into a verbal formulation that fits somewhere along a causal chain of federal powers is merely the beginning, not the end, of the constitutional inquiry. ... The terms `rationally related' and `rational basis' must be employed with care, particularly if either is to be used as a stand-alone test.''
\end{quote}
\noindent --- Justice Kennedy, concurring in the judgment (\href{https://www.law.cornell.edu/supct/html/08-1224.ZC.html}{560 U.S. at 150--51}).

\begin{quote}
``The Necessary and Proper Clause does not give Congress carte blanche. Although the term `necessary' does not mean `absolutely necessary' or indispensable, the term requires an `appropriate' link between a power conferred by the Constitution and the law enacted by Congress. ... Here, there is a substantial link to Congress' constitutional powers.''
\end{quote}
\noindent --- Justice Alito, concurring in the judgment (\href{https://www.law.cornell.edu/supct/html/08-1224.ZC1.html}{560 U.S. at 158}).

\begin{quote}
``For the reasons explained in Part III--A--2, infra, the enumerated power that justifies a particular defendant's criminal arrest or conviction cannot justify his subsequent civil detention under \S~4248.''
\end{quote}
\noindent --- Justice Thomas, dissenting (\href{https://www.law.cornell.edu/supct/html/08-1224.ZD.html}{560 U.S. at 167 n.5}).


\end{document}
